\chapter{Masa variable y ecuación de Tsiolkovski}\label{ch:25}
\objetivos{Derivar la ecuación de movimiento desde un sistema material; integrar el caso ideal; incluir fuerza externa y velocidad efectiva variable; distinguir presupuesto de delta-v y velocidad de una órbita.}
\section{Por qué no basta con derivar mv}
La segunda ley se aplica directamente a un sistema material de masa fija. Un vehículo que expulsa propulsante es un sistema abierto si se considera sólo la parte que permanece. Escribir $F=d(mv)/dt$ sin contabilizar el momento que cruza su frontera atribuye incorrectamente al vehículo el momento de la masa saliente.

Tomemos un intervalo $dt$ y un marco inercial. Al principio, vehículo y masa que va a expulsarse tienen masa conjunta $m$ y velocidad $v$. Al final, permanece una masa $m+dm$, con $dm<0$, a velocidad $v+dv$, y se ha expulsado masa $-dm$ con velocidad inercial $v-v_e$ a primer orden. La suma de momentos finales menos el inicial es
\begin{align}
 dP&=(m+dm)(v+dv)+(-dm)(v-v_e)-mv\\
 &=m\,dv+v_e\,dm+dm\,dv.
\end{align}
El producto de diferenciales es de segundo orden. Si no hay fuerzas externas distintas de la contribución de presión ya incorporada al empuje, resulta $m\,dv=-v_e\,dm$ en el caso de descarga adaptada. El signo permite acelerar hacia delante mientras la masa disminuye.
\section{Velocidad efectiva y fuerzas externas}
La presión de salida añade fuerza $(p_e-p_a)A_e$. Con caudal positivo $\dot m_p=-dm/dt$ y velocidad efectiva $c=F/\dot m_p$,
\begin{equation}\label{eq:rocketdiff}
 m\frac{dv}{dt}=-c\frac{dm}{dt}+F_{ext}.
\end{equation}
El empuje ya contiene el término de presión. Añadirlo otra vez a \eqref{eq:rocketdiff} lo duplicaría. Las fuerzas externas incluyen gravedad y resistencia según la frontera y el marco elegidos. La masa $m$ incluye estructura, carga útil, propulsante restante y todos los elementos que siguen unidos.

La aceleración inicial $F/m_0$ suele ser menor que la final si el empuje se mantiene y la masa baja. Una aceleración constante no se deduce de empuje constante. Tampoco la duración queda fijada sólo por $I_{sp}$: hace falta caudal o empuje y masa consumida.
\section{Integración ideal}
Si $F_{ext}=0$, el movimiento es rectilíneo, la velocidad efectiva es constante y no se desprende masa seca durante el intervalo,
\[
 dv=-c\frac{dm}{m}.
\]
Integrando entre $m_0,v_0$ y $m_f,v_f$,
\begin{equation}\label{eq:tsiolkovski}
 \boxed{\Delta v_{ideal}=v_f-v_0=c\ln\frac{m_0}{m_f}
 =g_0I_{sp}\ln\frac{m_0}{m_f}}.
\end{equation}
El logaritmo necesita un cociente adimensional. El resultado tiene unidades de velocidad. La ecuación es independiente del tiempo de expulsión bajo esas hipótesis; esa independencia desaparece cuando fuerzas externas o disponibilidad de potencia influyen en trayectoria.\cite{MIT32,JPL-EP}

La relación inversa es $m_0/m_f=\exp(\Delta v/c)$. La fracción de masa consumida es
\[
 \frac{m_p}{m_0}=1-e^{-\Delta v/c}.
\]
Para $\Delta v/c\ll1$, expandir la exponencial da $m_p/m_0\simeq\Delta v/c$. Esta aproximación pierde precisión cuando la fracción deja de ser pequeña. No sustituye al logaritmo en una etapa con gran consumo relativo.
\section{Sensibilidad y rendimientos decrecientes}
Con $m_f$ fijo,
\[
 \frac{\partial\Delta v}{\partial m_0}=\frac{c}{m_0},\qquad
 \frac{\partial^2\Delta v}{\partial m_0^2}=-\frac{c}{m_0^2}.
\]
Añadir masa inicial de propulsante produce una ganancia logarítmica y exige acelerar esa masa mientras permanece a bordo. Reducir masa seca puede tener un efecto intenso porque disminuye $m_f$ y, si cambia el sistema completo, también puede alterar otras masas.

A relación de masas fija, $\partial\Delta v/\partial c=\ln(m_0/m_f)$. Sin embargo, aumentar velocidad efectiva mediante una tecnología que añade mucha masa seca modifica la relación de masas. Comparar sólo el incremento de $I_{sp}$ sobreestima la ventaja si se mantiene ficticiamente el mismo presupuesto de masas.
\section{Velocidad efectiva variable}
Si $c$ cambia con estado, mezcla o ambiente,
\[
 \Delta v_{ideal}=\int_{m_f}^{m_0}\frac{c(m)}{m}\,dm.
\]
El promedio relevante para esta integral está ponderado por $dm/m$, no por masa consumida solamente ni por tiempo. El promedio que reproduce impulso total y el que reproduce delta-v no tienen por qué coincidir.

Para una descarga de potencia fija, $c$ puede fijarse aproximadamente y el empuje depender de potencia; si la potencia cambia, cambia duración y posiblemente régimen. Si se regula $c$, el caudal también cambia. La ecuación ideal sigue siendo un componente de la misión, pero no calcula por sí sola tiempo de transferencia.
\section{Delta-v y órbita}
La velocidad circular en un campo central es $v_c=\sqrt{\mu/r}$, obtenida al igualar aceleración centrípeta y gravitatoria. Es una propiedad de una órbita a radio $r$, no el delta-v propulsivo total desde una superficie. Un lanzamiento debe cambiar energía y momento angular, vencer efectos de gravedad durante el empuje y atravesar una atmósfera, con condiciones iniciales que incluyen rotación del planeta.

El presupuesto de delta-v suma maniobras y términos definidos para una misión. No es automáticamente el módulo de la velocidad final ni la diferencia escalar entre velocidades orbitales. En maniobras vectoriales, cambiar dirección puede consumir delta-v aun sin cambiar rapidez. Una órbita circular requiere dirección tangencial además de una rapidez apropiada.
\figura{c25}{Ganancia logarítmica del modelo ideal. No representa velocidad orbital alcanzada desde una superficie.}
\section{Caso desarrollado}
\begin{studiumworked}{M25. Capacidad de una etapa y masa requerida · INTERMEDIATE}
\textbf{Enunciado original.} Una etapa ideal tiene velocidad efectiva $2850\,\mathrm{m/s}$, masa inicial $1050\,\mathrm{kg}$ y final de combustión $365\,\mathrm{kg}$. Halle delta-v y fracción consumida. Para conservar esa masa final y exigir 1800 m/s, determine masa propulsante necesaria.

\textbf{Datos utilizados.} Velocidad efectiva: $2850\,\mathrm{m/s}$; Masa inicial: $1050\,\mathrm{kg}$; Masa final: $365\,\mathrm{kg}$; Objetivo ideal: $1800\,\mathrm{m/s}$.

\textbf{Incógnitas.} Capacidad ideal y masa inversa.

\textbf{Hipótesis.} Rectilíneo sin fuerzas externas; c constante. Masa final incluye todos los elementos unidos al terminar.

\textbf{Modelo físico.} Conservación de momento con masa variable.
\[m\,dv=-c\,dm,\quad \Delta v=c\ln(m_0/m_f)\]
\[m_{0,req}=m_f\exp(\Delta v_{req}/c)\]
\[m_{p,req}=m_f[\exp(\Delta v_{req}/c)-1]\]
\textbf{Despejes y sustituciones.}
\begin{enumerate}
\item Integrar $dv=-c\,dm/m$ entre las dos masas da $\Delta v=2850\ln(1050/365)=3011.447$ m/s.
\item La masa consumida inicial es $685\,\mathrm{kg}$ y su fracción $0.652381$. Para el objetivo distinto se despeja exponenciando el cociente adimensional $1800/2850$.
\item La nueva masa inicial es $686.4108\,\mathrm{kg}$; restar la masa final da $321.4108\,\mathrm{kg}$ de propulsante. Esa masa no debe sumarse a la masa inicial anterior: corresponde a un presupuesto nuevo con la misma masa final.
\end{enumerate}
\textbf{Resultados.} $\Delta v=3011.447$ m/s; fracción consumida $0.652381$; $m_{p,req}=321.4108$ kg.

\textbf{Interpretación.} El resultado es capacidad ideal, no velocidad orbital ni garantía de lanzamiento. Cambiar masa seca o pérdidas cambiaría el presupuesto inverso.

\textbf{Comprobación.} Sustituir $m_{0,req}/m_f=\exp(1800/c)$ recupera exactamente 1800 m/s. El logaritmo recibe un cociente positivo mayor que uno.

\textbf{Límites.} No hay separación de etapas ni fuerzas externas. La masa final se ha fijado como dato y no se recalcula estructura al cambiar propulsante.
\end{studiumworked}

\section{Ejercicios propuestos}
\begin{enumerate}
\item Repita la derivación con $dm<0$ y compruebe el signo del incremento de velocidad.
\item Calcule fracción consumida para $\Delta v/c=0.2$ y compare con la aproximación lineal.
\item Obtenga la masa de propulsante requerida a partir de una masa final dada.
\item Explique por qué un delta-v ideal igual a velocidad circular no demuestra que un vehículo alcance esa órbita desde una superficie.
\end{enumerate}
\resumen{La masa saliente transporta momento y debe contabilizarse. Tsiolkovski integra un modelo sin fuerzas externas y con velocidad efectiva constante. Su delta-v es una capacidad ideal de maniobra, distinta de velocidad orbital, tiempo de vuelo y prestación real de lanzamiento.}


\studiumchaptercheck
